\documentclass[11pt]{article}
\usepackage[margin=1in]{geometry}
\usepackage{amsmath,amssymb}
\title{有限情景下的投资与运营模型}
\author{}
\date{}
\usepackage[UTF8]{ctex}
\begin{document}
\maketitle

\section{集合与索引}
让$\mathcal S$表示有限的损害场景集合。配电网络为$(N,E)$，通信网络为$(\widetilde N,\widetilde E)$，其中包括虚拟节点$\widetilde i_d$。用$\mathcal E=E\cup\widetilde E$表示集成边缘集。每条边的端点顺序分别为$(e_h,e_t)$、$\delta(e)=\{e_h,e_t\}$，而$\delta(i)$则表示在$i\in N$时发生的通信事件。在由$e=(i,j)$、$i=e_h$和$j=e_t$索引的表达式中，通信事件的总和范围为$\widetilde E$。

现有的和候选的边缘集合分别是$E_x^0$和$E_x^n$，还有$E_x^0\cup E_x^n=\mathcal E$。切换集合则包括$E_t^0,E_t^n\subseteq E$和$E_t=E_t^0\cup E_t^n$。符合加固条件的集合是$E_h\subseteq\mathcal E$。现有的和候选的生成器组合起来就是$U=U^0\cup U^n$。这些现有/候选的分区是不相交的。每个生成器$l$都位于总线$i(l)$和$U_i=\{l\in U:i(l)=i\}$上。候选的分配线路具有现有的切换可用性，具体信息由$E_x^n\cap E\subseteq E_t^0$所指定。

让$E_V\subseteq E$表示变压器集合，$L\subseteq N$表示关键负载总线，而$D_s\subseteq\mathcal E$则表示在场景$s$中受损的边。让$\mathcal C$表示所提供的有限数量的分配周期集合，而$E_{\mathcal C}=\bigcup_{C\in\mathcal C}C$则表示其他元素。拓扑约束中的端点匹配包括并行边。

设$P=\{a,b,c\}$，其包含总线、边和生成器阶段集合$P_i,P_e,P_l\subseteq P$。变压器阶段集合是非空的。在阶段索引表达式中，缺失的阶段用零来替代。所有阶段的投影和扩展都遵循这一约定。

通信控制点集为$\widetilde N_c\subseteq\widetilde N$，其中$\widetilde i(l)\in\widetilde N_c$对应$l\in U$，而$\widetilde i(e)\in\widetilde N_c$对应$e\in E_t$。

\section{参数}
投资成本分别为：$c_e^x$用于$e\in E_x^n$，$c_e^t$用于$e\in E_t^n$，$c_e^h$用于$e\in E_h$，$c_l^u$用于$l\in U^n$。实际和反应性发电机的产能分别为：$\bar g_{l,p}^k,\bar g_{l,q}^k$用于$l\in U$，$k\in P_l$。需求分别为：$D_{i,p}^k,D_{i,q}^k$用于$i\in N$，$k\in P_i$。关键和服务的总比例分别为：$\eta_c,\eta_t$。

对于$e\in E$和$k\in P_e$，$T_e^k$是热极限参数（在源代码中也表示为$T_{e,k}$）。变压器共享系数为$\underline\beta_e,\overline\beta_e$，适用于$e\in E_V$。电压幅值范围分别为$\underline V_i^k,\overline V_i^k$，适用于$i\in N$和$k\in P_i$。每个阻抗矩阵都表示为$Z_e=R_e+\mathrm i X_e$，其中$R_e,X_e\in\mathbb R^{P_e\times P_e}$和$\mathrm i^2=-1$是其中的具体数值。

使用平衡的相位向量$u_a=1$、$u_b=\alpha$、$u_c=\alpha^2$，其中$\alpha=\exp(-2\pi\mathrm i/3)$，并对其进行定义。
\[
 \Gamma_{kk'}=u_k/u_{k'},\quad
 C_{e,kk'}=2\Gamma_{kk'}\overline{(Z_e)_{kk'}}.
\]
物理上的实际平方电压降为$\Delta_e^{s,k}=\sum_{k'\in P_e}
[\operatorname{Re}(C_{e,kk'})p_{e,i}^{s,k'}
-\operatorname{Im}(C_{e,kk'})q_{e,i}^{s,k'}]$。这是$S_{e,i}Z_e^H+Z_eS_{e,i}^H$在平衡电压、零损耗近似条件下的对角线值（参见第13-14页，方程(7)-(10)）。

让$K$成为一个固定的正整数多边形参数，该参数会与实例一起提供（所报告的实现方式使用的是$K=28$）。请预先计算好这个值。
\[
a_n=\sin\frac{2n\pi}{K}-\sin\frac{2(n-1)\pi}{K},\qquad
b_n=\cos\frac{2n\pi}{K}-\cos\frac{2(n-1)\pi}{K},\qquad
r_K=\sin\frac{2\pi}{K},\quad n=1,\ldots,K.
\]
使用源代码的有限限制条件。
\[
M_{\rm comm}=|\widetilde N_c|,\qquad
M_{\rm flow}^k=\sum_{i\in N}D_{i,p}^k\quad(k\in P).
\]
所有集合和参数数组都是有限输入数据。

\section{变量和域}
第一阶段二进制投资向量和匹配成本向量是
\[
w=(x_{E_x^n},t_{E_t^n},h_{E_h},u_{U^n})\in\{0,1\}^m,
\qquad c=(c^x_{E_x^n},c^t_{E_t^n},c^h_{E_h},c^u_{U^n}),
\]
其中$m=|E_x^n|+|E_t^n|+|E_h|+|U^n|$，采用固定的共同排序方式。

对于每一个$s\in\mathcal S$，场景中的二进制变量为：
\[
x_e^s,t_e^s,h_e^s\in\{0,1\}\ (e\in\mathcal E),\qquad
u_l^s\in\{0,1\}\ (l\in U),
\]
\[
z_e^s,b_e^s,b_e^{0,s}\in\{0,1\}\ (e\in E),\qquad
y_e^s\in\{0,1\}\ (e\in E_{\mathcal C}).
\]
这里$x^s$表示边可用性，$t^s$表示开关使用情况，$h^s$表示加固措施，$u^s$表示发电机的可用性，$z^s$表示边活动情况，$b^s,b^{0,s}$表示流向指示符，而$y^s$则表示循环指示符。定义向量缩写。
\[
w^s=(x^s_{E_x^n},t^s_{E_t^n},h^s_{E_h},u^s_{U^n}).
\]
真正的连续变量，每个情景下都是分开的。
\begin{align*}
&p_{e,i}^{s,k},q_{e,i}^{s,k} &&(e\in E,\ i\in\delta(e),\ k\in P_e),\\
&g_{l,p}^{s,k},g_{l,q}^{s,k} &&(l\in U,\ k\in P_l),\\
&d_{i,p}^{s,k},d_{i,q}^{s,k},v_i^{s,k} &&(i\in N,\ k\in P_i),\\
&f_e^s &&(e\in\widetilde E),\\
&\gamma_i^s &&(i\in\widetilde N_c).
\end{align*}
它们代表了端点实值/反应性流、生成、传递的负载、平方后的电压幅度、人工通信流以及控制点的连通性。连续变量为实值，其上下限在下面的约束条件中有所规定。

\section{目标与约束条件}
该模型能够最小化
\[
\min\ c^\top w.
\]
投资链接是：
\[
w\geq w^s\qquad(s\in\mathcal S).
\]
本节剩余部分中的每个约束都是针对每一个$s\in\mathcal S$而设定的。

\subsection{设备固定装置}
\begin{align*}
x_e^s&=1 &&(e\in E_x^0\setminus D_s),\\
t_e^s&=0 &&(e\in\mathcal E\setminus E_t),\\
h_e^s&=0 &&(e\in\mathcal E\setminus E_h),\\
u_l^s&=1 &&(l\in U^0).
\end{align*}

\subsection{功率平衡与方向限制}
对于$i\in N$和$k\in P_i$，
\begin{align*}
\sum_{l\in U_i}g_{l,p}^{s,k}-d_{i,p}^{s,k}
 &=\sum_{e\in\delta(i)}p_{e,i}^{s,k},\\
\sum_{l\in U_i}g_{l,q}^{s,k}-d_{i,q}^{s,k}
 &=\sum_{e\in\delta(i)}q_{e,i}^{s,k}.
\end{align*}
对于每一条有序的边$e=(i,j)\in E$和$k\in P_e$，
\begin{align*}
-T_e^k b_e^s&\leq p_{e,i}^{s,k}\leq T_e^k(1-b_e^s),\\
-T_e^k b_e^{0,s}&\leq q_{e,i}^{s,k}\leq T_e^k(1-b_e^{0,s}).
\end{align*}

\subsection{无损端点保留}
端点流是从它们自己的总线上以正数值流出。因此，对于每一种情况来说，物理线路$e=(i,j)$和相位$k\in P_e$都是如此。
\[
 p_{e,i}^{s,k}+p_{e,j}^{s,k}=0,\qquad
 q_{e,i}^{s,k}+q_{e,j}^{s,k}=0.
\]
分布式自环是无效的。现在对节点实功率平衡进行求和，可证明总发电量等于总需求；独立的端点变量无法产生电力。

\subsection{相位共享}
对于有序的变换器$e=(i,j)\in E_V$，请使用以下缩写形式表示：
\[
H_{e,p}^s=\sum_{r\in P_e}p_{e,i}^{s,r},\qquad
H_{e,q}^s=\sum_{r\in P_e}q_{e,i}^{s,r},\qquad d_e=|P_e|.
\]
对于每一个$k\in P_e$，都施加扩展的二次不等式。
\begin{align*}
\underline\beta_e H_{e,p}^s+
 (\overline\beta_e-\underline\beta_e)b_e^sH_{e,p}^s
 &\leq d_e p_{e,i}^{s,k},\\
d_e p_{e,i}^{s,k}
 &\leq\overline\beta_e H_{e,p}^s+
 (\underline\beta_e-\overline\beta_e)b_e^sH_{e,p}^s,\\
\underline\beta_e H_{e,q}^s+
 (\overline\beta_e-\underline\beta_e)b_e^{0,s}H_{e,q}^s
 &\leq d_e q_{e,i}^{s,k},\\
d_e q_{e,i}^{s,k}
 &\leq\overline\beta_e H_{e,q}^s+
 (\underline\beta_e-\overline\beta_e)b_e^{0,s}H_{e,q}^s.
\end{align*}

\subsection{生成、加载和服务}
\begin{align*}
0\leq g_{l,p}^{s,k}&\leq\bar g_{l,p}^k u_l^s,&
g_{l,q}^{s,k}&\leq\bar g_{l,q}^k u_l^s
&& (l\in U,\ k\in P_l),\\
0\leq d_{i,p}^{s,k}&\leq D_{i,p}^k,&
0\leq d_{i,q}^{s,k}&\leq D_{i,q}^k
&& (i\in N,\ k\in P_i).
\end{align*}
对于每一个$k\in P$和$a\in\{p,q\}$，
\[
\sum_{i\in L}d_{i,a}^{s,k}\geq\eta_c\sum_{i\in L}D_{i,a}^k,
\qquad
\sum_{i\in N}d_{i,a}^{s,k}\geq\eta_t\sum_{i\in N}D_{i,a}^k.
\]

\subsection{通信}
\begin{align*}
\sum_{e\in\widetilde E:e_t=i}f_e^s-
\sum_{e\in\widetilde E:e_h=i}f_e^s
 &=\gamma_i^s &&(i\in\widetilde N_c),\\
\sum_{e\in\widetilde E:e_t=i}f_e^s-
\sum_{e\in\widetilde E:e_h=i}f_e^s
 &=0 &&(i\in\widetilde N\setminus(\widetilde N_c\cup\{\widetilde i_d\})),\\
-M_{\rm comm}x_e^s\leq f_e^s
 &\leq M_{\rm comm}x_e^s &&(e\in\widetilde E),\\
0\leq g_{l,p}^{s,k}
 &\leq\bar g_{l,p}^k\gamma_{\widetilde i(l)}^s
 &&(l\in U,\ k\in P_l),\\
g_{l,q}^{s,k}
 &\leq\bar g_{l,q}^k\gamma_{\widetilde i(l)}^s
 &&(l\in U,\ k\in P_l),\\
t_e^s&\leq\gamma_{\widetilde i(e)}^s &&(e\in E_t),\\
0\leq\gamma_i^s&\leq1 &&(i\in\widetilde N_c).
\end{align*}

\subsection{可用性与拓扑结构}
\begin{align*}
x_e^s&\geq t_e^s &&(e\in\mathcal E),\\
z_e^s&=x_e^s-t_e^s &&(e\in E),\\
x_e^s&=h_e^s &&(e\in D_s),\\
\sum_{e\in C}y_e^s&\leq |C|-1 &&(C\in\mathcal C),\\
z_{\hat e}^s&\leq y_e^s
 &&(C\in\mathcal C,\ e\in C,\ \hat e\in E:\delta(\hat e)=\delta(e)).
\end{align*}

\subsection{活动线路上的电压}
对于每一个$e=(i,j)\in E$、$k\in P_e$，进行定义。
\[
 M_{e,k}^{v}=\max\{(\bar V_i^k)^2-(\underline V_j^k)^2,
                    (\bar V_j^k)^2-(\underline V_i^k)^2\}.
\]
强制实施/执行
\[
-M_{e,k}^{v}(1-z_e^s)
\le v_i^{s,k}-v_j^{s,k}-\Delta_e^{s,k}
\le M_{e,k}^{v}(1-z_e^s),
\qquad
(\underline V_i^k)^2\le v_i^{s,k}\le(\bar V_i^k)^2.
\]
在活动边界条件下，无效线路端点的功率已经为零，因此$M^v$允许在线路断开时存在任何允许的端点电压组合。之前的无条件等式以及多余的虚部等式都被移除了。源代码中明确说明，欧姆定律仅适用于可用的线路（见第10页，(2a)；第14页，紧随(10)之后）。

\subsection{多边形容量和活动限制}
对于$e\in E$、$i\in\delta(e)$和$k\in P_e$，
\begin{align*}
a_n p_{e,i}^{s,k}-b_n q_{e,i}^{s,k}
 &\leq r_K T_e^k &&(n=1,\ldots,K),\\
-M_{\rm flow}^k z_e^s\leq p_{e,i}^{s,k}
 &\leq M_{\rm flow}^k z_e^s,\\
-M_{\rm flow}^k z_e^s\leq q_{e,i}^{s,k}
 &\leq M_{\rm flow}^k z_e^s.
\end{align*}

\paragraph{明确的物理解释。}
这是一个平衡电压、无损耗的三相近似模型，其中包含了源头的有限损伤情景、投资链接、控制网络、辐射拓扑结构、负载服务比例以及内部多边形的热边界条件。二进制连续相位共享产品被完全保留下来，同时采用了全局MIQCP处理方式。端点流量采用外向流符号约定。实际电压降的关系可以通过(7)--(8)式得出。
\end{document}
